% State of the Art section. Paste into your own thesis, or \input it.
%
% Objective 4's segmentation half: one paragraph per method bransby2026imagecasx benchmarks in
% its Table 2/3, each cited to its own original paper, then that benchmark comparison itself,
% closing on why CAS-Net is the method this thesis's framework builds on (CLAUDE.md "The
% framework half"). The topological-description half of objective 4 is the rest of this chapter.
%
% Citations for the six benchmarked methods (CAS-Net, the 3D-FFR-UNet, TotalSegmentator, ADE-HTL,
% nnU-Net, Swin UNETR) were pulled 2026-09-07 from bransby2026imagecasx's own numbered reference
% list (its [11], [12], [29], [31], [32], [33]) and verified against Crossref -- see the
% "six methods" block in thesis/refs.bib. None of the six is [VERIFIED] (read in full); only
% bransby2026imagecasx itself, whose Table 2/3 this section reports, carries that status here.
% All quantitative results below are mean $\pm$ standard deviation over the same 160 test cases,
% as reported in bransby2026imagecasx Tables 2 and 3.

\section{Segmentation methods}
\label{sec:segmentation}

A coronary tree cannot be built until the lumen is segmented, so the choice of segmentation
method is the first decision this thesis's framework makes.
Objective 5 is to adopt an established method rather than design a new one, and ImageCAS-X's own
technical validation makes that choice a benchmarked one rather than a guess: six recent
segmentation methods, standardized to one training pipeline and scored on the same 160 held-out
test cases, compared against each other, against inter-observer agreement, and against the label
set the original ImageCAS dataset shipped \cite{bransby2026imagecasx}.

Three measures carry that comparison.
The \textbf{Dice similarity coefficient (DSC)} is the fraction of predicted and true lumen voxels
that overlap.
The 95th-percentile \textbf{Hausdorff distance (HD95)} is the surface error a large local mistake
produces, in millimeters, taken at the 95th percentile so a handful of stray voxels cannot dominate
it.
The \textbf{average symmetric surface distance (ASSD)} is the mean of that same surface error
instead of its tail.
Two further measures score the extracted centerline rather than the lumen volume: clDice, a
Dice-like overlap between predicted and true centerlines \cite{shit2021cldice}, and the Betti
number error, which counts breaks and false loops in the predicted tree and so penalizes a
topological mistake that a voxel-overlap score can miss entirely
\cite{bransby2026imagecasx}.

Two of the six methods classify patches sampled at random from the full volume, so neither has
direct access to the tree's global shape while segmenting a given patch.
Song et al.'s 3D-FFR-UNet stacks two U-Nets, one to produce a coarse prediction and a second to
rectify it, reaching DSC $84.9 \pm 5.5$ \cite{song2022ffrunet} --- despite the name, "FFR" here is
Feature-Fusion-and-Rectification, not fractional flow reserve (Section~\ref{sec:ffr}), and the
network does not estimate a pressure ratio.
Dong et al.'s CAS-Net instead attends across three input scales at once, so a patch is classified
using both fine local detail and coarser surrounding context, and it is the highest-scoring of all
six methods on DSC and ASSD \cite{dong2023casnet}.

Zeng et al.'s ImageCAS baseline, released with the original ImageCAS dataset
(Section~\ref{sec:dataset-cohort}), segments coarse-to-fine instead: a coarse whole-volume pass
first, then a second network that classifies patches centered on that coarse prediction, giving
distal vessels a targeted second pass rather than an equal share of randomly sampled patches
\cite{zeng2023imagecas}.
It reaches DSC $87.9 \pm 2.9$, in the middle of the six.

Zhang et al.'s ADE-HTL starts from the centerline rather than the volume: it locates bifurcation
points explicitly and enforces the tree's connectivity as part of the network, instead of leaving
connectivity to emerge from voxels classified independently \cite{zhang2024adehtl}.
That choice makes ADE-HTL the closest competitor to CAS-Net overall, and on the two metrics that
describe the segmentation as a topological object rather than a volume it is the more accurate of
the two: Betti number error $1.5 \pm 1.5$ against CAS-Net's $1.9 \pm 1.5$, a significant difference
($p<0.001$) and CAS-Net's only significant loss to any other method
\cite{bransby2026imagecasx,zhang2024adehtl}.
On DSC the result reverses and is CAS-Net's largest significant margin over any method,
$91.2 \pm 2.8$ against ADE-HTL's $87.7 \pm 2.8$ ($p<0.001$); on HD95, clDice, ASSD and centerline
HD95 the two methods are not significantly different \cite{bransby2026imagecasx}.

Two of the six methods are not coronary-specific at all.
nnU-Net configures its own architecture, preprocessing and training schedule from whatever dataset
it is given, rather than being hand-tuned for one anatomical target \cite{isensee2021nnunet}, and
Swin UNETR replaces a U-Net's convolutional encoder with a transformer, built originally for brain
tumors on magnetic resonance imaging \cite{hatamizadeh2022swinunetr}.
Retrained here on the coronary lumen, both land in the middle of the six: nnU-Net reaches DSC
$89.8 \pm 3.2$, and Swin UNETR $87.9 \pm 2.7$.
Replacing nnU-Net's loss with clDice, which penalizes a broken centerline directly rather than only
voxel overlap, raises its DSC to $90.0 \pm 3.6$ without closing the gap to CAS-Net
\cite{shit2021cldice,bransby2026imagecasx}.

TotalSegmentator is the only one of the six not trained on the coronary lumen at all: a general
anatomical segmentation network for 104 structures, applied here zero-shot to a structure outside
that set \cite{wasserthal2023totalsegmentator}.
It scores lowest of the six on every measure reported, DSC $70.5 \pm 6.2$ against CAS-Net's
$91.2 \pm 2.8$ ($p<0.001$), which the benchmark's authors attribute to a possible distribution
shift between its training data and coronary computed tomography angiography rather than to a
general limit on pretrained segmentation networks \cite{bransby2026imagecasx}.

CAS-Net's own numbers, in full: DSC $91.2 \pm 2.8$, HD95 $2.99 \pm 3.47$ mm, Betti number error
$1.9 \pm 1.5$, clDice $93.3 \pm 3.2$, ASSD $0.73 \pm 0.36$ mm, centerline HD95
$5.75 \pm 4.98$ mm \cite{bransby2026imagecasx,dong2023casnet}.
That result needs a ceiling and a floor.
The ceiling is inter-observer agreement: two independent analysts, each starting from the same
automatically generated centerlines and the same initial network prediction, agreed with each
other at DSC $92.8 \pm 3.1$ on the same 160 cases \cite{bransby2026imagecasx}.
CAS-Net does not reach it — the gap is significant on DSC, HD95 and Betti number error
($p<0.001$ each) — and because both analysts edited the same starting point rather than annotating
independently from nothing, this figure is an upper bound on demonstrated agreement, not a neutral
one \cite{bransby2026imagecasx}.
The floor is the label set the original ImageCAS release shipped: scored against the same
relabeled test cases, those labels reach only DSC $41.8 \pm 6.7$
\cite{bransby2026imagecasx,zeng2023imagecas} --- the gap ImageCAS-X's own annotation effort
(Section~\ref{sec:dataset-cohort}) exists to close.

Accuracy is not the only cost that matters for a framework meant to run on thousands of scans.
CAS-Net is also the cheapest of the six: 6.8 million trainable parameters against 10.7--153.9
million for the rest, and 15.5 seconds per scan against up to 102.3 seconds for the two nnU-Net
variants and 2100 seconds, 35 minutes, for a human analyst \cite{bransby2026imagecasx}.
ADE-HTL's one significant advantage over CAS-Net, on Betti number error, comes with a cost of its
own: this framework's implementation of it additionally requires an anatomical prior from
TotalSegmentator's heart-chambers model, one more dependency than CAS-Net needs, and is not yet
fully wired up to run.
Highest accuracy on every metric that is not tied, at the lowest computational cost and the fewest
dependencies, is why this thesis's framework builds its segmentation stage on CAS-Net rather than
on any of the other five.
